\section{链接自能贡献与紫外发散}

最简单的图对应于规范链接的
自能修正（见图 \ref{linkself}）。其计算与夸克双定域算符情形相同
（例如见 Ref.~\cite{Radyushkin:2017lvu}）。
单圈水平上只需把色因子改为 $C_F \to C_A$。于是我们有
\begin{align}
\Gamma_ \Sigma (z) = & ({i} g)^2\,C_A \,\frac12 \, \int_0^1 \dd t_1 \, \int_0^1 \dd t_2 \,
z^\mu z^\nu \, D_{\mu \nu} ^c [ z (t_2-t_1) ] \ ,
\label{self}
\end{align}
其中 $D_{\mu \nu} ^c (z) =g_{\mu \nu} /4\pi^2 z^2$
是坐标表象中的费曼规范胶子传播子。
对链接参数 $t_1, t_2$ 的积分
\begin{align}
\int_0^1 \dd t_1 \, \int_0^1 \frac{\dd t_2}{(t_2-t_1)^2} \,
\label{t1t2}
\end{align}
在 $t_1 \sim t_2$ 时发散，即当
胶子传播子的端点 $t_1z$ 与 $t_2 z$ 彼此靠近时发散。因此，
可以怀疑这一发散具有紫外起源。
为确认确实如此，我们在 UV 区域采用维数正规化（DR）\cite{tHooft:1972tcz}，
切换到 $d$ 维。
这样，坐标空间中的胶子传播子获得额外因子 $(-z^2)^{2-d/2}$。
这在式 (\ref{t1t2}) 中引入额外因子 $(t_2-t_1)^{2-d/2}$，使该积分在足够小的 $d$ 下收敛。

\begin{figure}[t]
\centerline{\includegraphics[width=2in]{images/linkgluon}}
\vspace{-8mm}
\caption{规范链接的自能型修正。
\label{linkself}}
\end{figure}

为保持规范不变性，我们的计算使用无质量胶子与维数正规化。
然而，就链接自能图而言，使用 DR
（它基本上只是一种数学技巧）在若干地方颇具误导性。

相关的微妙之处可以通过对坐标表象中的胶子传播子采用 Polyakov 规定
\mbox{$1/z^2 \to 1/(z^2-a^2)$} \cite{Polyakov:1980ca}
（另见文献 \cite{Chen:2016fxx,Radyushkin:2017lvu}）来说明。
它使胶子传播子在 $-z^2 \lesssim a^2$ 区间变软，
并消除其在 $z^2=0$ 处的奇异性。就此而言，
它类似于有限格距产生的 UV 正规化，并给出
\begin{align}
\Gamma_ \Sigma (z,a) = & - g^2\,C_A \,\frac{z^2}{8 \pi^2}
\, \int_0^1 \dd t_1 \, \int_0^1 \, \frac{\dd t_2}{ z^2 (t_2-t_1)^2 - a^2} \ .
\label{selfa}
\end{align}
正规化后的积分在光锥 \mbox{$z^2=0$} 上为零，并对类空 $z$ 收敛。
取 $z=z_3$ 并完成积分得 \cite{Chen:2016fxx,Radyushkin:2017lvu}
\begin{align}
\Gamma_\Sigma (z_3,a) = -\,C_A \,\frac{\alpha_s}{2 \pi}
&
\left [
\,
2 \frac{z_3}{a} \, \tan
^{-1}\left(\frac{z_3}{a}\right) - \ln
\left(1+ \frac{z_3^2}{a^2}\right) \right ] \ .
\label{selfex}
\end{align}
结果含有 DR 所遗漏的线性 $\sim 1/a$ 发散。此外，
固定 $a$ 而令 $z_3$ 变小时，其行为如 $z_3^2/a^2$，即 $\Gamma_\Sigma (z,a)$ 在 $z_3=0$ 时为零，
正如预期：$z_3=0$ 时不存在链接。
它在光锥 $z^2=0$ 上也为零。

$\Gamma_\Sigma (z_3=0,a) =0$ 这一事实意味着：
对固定的 $a$，
此项对
$G_{\mu \alpha} (z) \, [z,0]\, G_{\lambda \beta} (0)$
算符的定域极限（例如对能量动量张量（EMT））不产生修正。
由于 EMT 的矩阵元给出
胶子携带的强子动量份额，
链接自能修正不改变这一份额。在没有
胶子--夸克跃迁时这是自然的现象。

然而，若在式 (\ref{selfex}) 中形式地对固定 $z_3$ 取 $a\to 0$ 极限，
则 $\ln \left(1+{z_3^2}/{a^2}\right)$ 变为在 $z_3=0$ 奇异的表达式
$\ln \left({z_3^2}/{a^2}\right)$。
类似地，使用 DR 会遇到正比于
\begin{align}
(-z^2\mu_{\rm UV}^2)^{\euv}/\euv= 1/\euv + \ln (-z^2 \mu_{\rm UV}^2) + \ldots \, ,
\end{align}
的结果，其中 $\mu_{\rm UV}$ 是伴随此维数 UV 正规化的标度。
同样，初始表达式在 $z^2=0$ 处为零，但通过减除 $1/\euv$ 极进行重正化后，
人们似乎会得出结论：除 UV 发散外，该图还含光锥 $z^2=0$ 上的奇异性。

出于这一原因，在我们的 DR 结果中，我们将明确区分
由 UV 奇异项诱导的 \mbox{$z^2$ 依赖}
（它们实际上在光锥上消失）
与出现在 DGLAP 演化
对数 $\ln (-z^2 \mu_{\rm IR}^2)$ 中的 \mbox{$z^2$ 依赖}，这里 $\mu_{\rm IR}$ 是与共线奇异性的
DR 正规化相关联的标度。

主要区别在于：如果不用 DR，而用物理红外截断 $\Lambda$（如非零胶子虚度或胶子质量）
来正规化共线奇异性，
则正比于修正贝塞尔函数 $K_0 (\sqrt{-z^2 \Lambda^2})$ 的单圈结果在 $z^2=0$ 处仍是奇异的，
这与 UV 诱导的对数 $\ln (1-z^2/a^2)$ 不同。

就链接自能图而言，只存在 UV 奇异性。
它对 $G_{\mu \alpha}(z)G_{\lambda \beta}(0)$ 算符的修正由下式给出：
\begin{align}
&
-{g^2N_c \over 4\pi^2[(-z^2 \mu_{\rm UV}^2 +i\epsilon)]^{{d\over 2}-2}} {\Gamma\big({d/ 2}-1\big) \over (3-d)(4-d)}G_{\mu\alpha}(z)G_{\lambda \beta }(0) \ ,
\label{selfAD}
\end{align}
其中因子 $1/(3-d)(4-d)$ 来自
对胶子传播子 \mbox{$D^c (t_1z - t_2 z)$} 作 DR 产生的
积分
\begin{align}
&
\!\int_0^1\! \dd t_1 \!\int_0^{t_1}\! \dd t_2\, (t_1-t_2)^{2-d} ={1 \over (3-d)(4-d)}
\nonumber
\end{align}
$d=3$（$d=4$）处的极点对应于式 (\ref{selfex}) 中的线性（对数）UV 发散。
